\begin{definition} \label{an1:def:metric-space}
  A function $d: X \times X \to [0, \infty)$ is called a \textbf{metric} on $X$ if
  \begin{enumerate}[label=\textbf{(\Alph*)}]
    \item $\forall p, q \in X, d(p, q) = 0 \iff p = q$
    \item $\forall p, q \in X, d(p, q) = d(q, p)$
    \item $\forall p, q, r \in X, d(p, r) \leq d(p, q) + d(q, r)$
  \end{enumerate}
  \textbf{(C)} is called the \textbf{triangle inequality}. The structure $(X, d)$ is called a \textbf{metric space}.

\end{definition}

\begin{theorem} \label{an1:thm:standard-metric-on-Rn}
  Suppose $n \in \N$. Define $d: \R^{n} \times \R^{n} \to [0, \infty)$ by
  \[
    \forall \bm{x}, \bm{y} \in \R^{n}, d(\bm{x}, \bm{y}) \coloneqq \norm*{\bm{x} - \bm{y}}
  \]
  Then $d$ is a metric on $\R^{n}$.
\end{theorem}

\begin{proof}
  By \cref{an1:def:metric-space},
  \begin{itemize}
    \item Suppose $\bm{x}, \bm{y} \in \R^{n}$. By \cref{an1:def:inner-product-norm},
          \[
            \norm*{\bm{x} - \bm{y}} = 0 \iff \forall k \in [1, n]_{\N}, x_{k} = y_{k} \iff \bm{x} = \bm{y}
          \]
    \item Suppose $\bm{x}, \bm{y} \in \R^{n}$. By \cref{an1:def:inner-product-norm},
          \[
            \norm*{\bm{x} - \bm{y}} = \sqrt{\sum_{k = 1}^{n} (x_{k} - y_{k})^{2}}
            = \sqrt{\sum_{k = 1}^{n} (y_{k} - x_{k})^{2}} = \norm*{\bm{y} - \bm{x}}
          \]
    \item Suppose $\bm{x}, \bm{y}, \bm{z} \in \R^{n}$. By \cref{an1:thm:properties-of-inner-product-norm},
          \[
            \norm*{\bm{x} - \bm{z}} = \norm*{(\bm{x} - \bm{y}) + (\bm{y} - \bm{z})}
            \leq \norm*{\bm{x} - \bm{y}} + \norm*{\bm{y} - \bm{z}}
          \]
  \end{itemize}
\end{proof}

\begin{remark}
  Throughout this textbook, \textbf{euclidean spaces} use the \textbf{standard metric}
  defined in \cref{an1:thm:standard-metric-on-Rn}, unless otherwise specified.
\end{remark}

\begin{definition} \label{an1:def:neighborhood-open-ball}
  Suppose
  \begin{itemize}
    \item $p \in \abs*{(X, d)}$
    \item $r \in \R^{+}$
  \end{itemize}
  Then
  \begin{itemize}
    \item The \textbf{neighborhood} of $p$ with radius $r$ is defined by
          \[
            N_{r}(p) \coloneqq \setbuilder*{q \in X}{d(p, q) < r}
          \]
    \item The \textbf{deleted neighborhood} of $p$ with radius $r$ is defined by
          \[
            \widetilde{N}_{r}(p) \coloneqq N_{r}(p) \setminus \set*{p}
            = \setbuilder*{q \in X}{0 < d(p, q) < r}
          \]
  \end{itemize}
\end{definition}

\begin{definition} \label{an1:def:limit-interior-isolated-points}
  Suppose
  \begin{itemize}
    \item $E \subseteq \abs*{(X, d)}$
    \item $p \in X$
  \end{itemize}
  Then
  \begin{itemize}
    \item $p$ is called an \textbf{limit point} of $E$ if
          \[
            \forall r \in \R^{+}, E \cap \widetilde{N}_{r}(p) \neq \vnot
          \]
          The set of all limit points of $E$ is denoted by
          \[
            E' \coloneqq \setbuilder*{p \in X}{\forall r \in \R^{+}, E \cap \widetilde{N}_{r}(p) \neq \vnot}
          \]
    \item $p$ is called an \textbf{interior point} of $E$ if
          \[
            \exists r \in \R^{+}, N_{r}(p) \subseteq E
          \]
          The set of all interior points of $E$ is denoted by
          \[
            E^{\circ} \coloneqq \setbuilder*{p \in X}{\exists r \in \R^{+}, N_{r}(p) \subseteq E}
          \]
    \item $p$ is called a \textbf{isolated point} of $E$ if
          \[
            \exists r \in \R^{+}, E \cap N_{r}(p) = \set*{p}
          \]
          The set of all isolated points of $E$ is denoted by
          \[
            E^{i} \coloneqq \setbuilder*{p \in X}{\exists r \in \R^{+}, E \cap N_{r}(p) = \set*{p}}
          \]
  \end{itemize}
\end{definition}

\begin{definition} \label{an1:def:open-closed-sets}
  Suppose $E \subseteq \abs*{(X, d)}$. Then
  \begin{itemize}
    \item $E$ is said to be \textbf{open} if
          \[
            \forall p \in E, p \in E^{\circ}
          \]
          The set of all open sets in $X$ is denoted by
          \[
            \mathcal{O}(X) \coloneqq \setbuilder*{E \subseteq X}{\forall p \in E, p \in E^{\circ}}
          \]
    \item $E$ is said to be \textbf{closed} if
          \[
            E' \subseteq E
          \]
          The set of all closed sets in $X$ is denoted by
          \[
            \mathcal{C}(X) \coloneqq \setbuilder*{E \subseteq X}{E' \subseteq E}
          \]
  \end{itemize}
\end{definition}

\begin{definition} \label{an1:def:closure}
  Suppose $E \subseteq \abs*{(X, d)}$. The \textbf{closure} $\overline{E}$ of $E$ is defined by
  \[
    \overline{E} \coloneqq E \cup E'
  \]
\end{definition}

\begin{definition} \label{an1:def:bounded-perfect-dense-sets}
  Suppose $E \subseteq \abs*{(X, d)}$. Then
  \begin{itemize}
    \item $E$ is said to be \textbf{bounded} if
          \[
            X \neq \vnot \implies \exists p \in X, \exists r \in \R^{+}, E \subseteq N_{r}(p)
          \]
          The set of all bounded sets in $X$ is denoted by
          \[
            \mathcal{B}(X) \coloneqq \setbuilder*{E \subseteq X}{
              X \neq \vnot \implies \exists p \in X, \exists r \in \R^{+}, E \subseteq N_{r}(p)
            }
          \]
    \item $E$ is said to be \textbf{perfect} if
          \[
            E = E'
          \]
    \item $E$ is said to be \textbf{dense} if
          \[
            \overline{E} = X
          \]
  \end{itemize}
\end{definition}

\begin{theorem} \label{an1:thm:neighborhood-is-open}
  \[
    \forall p \in \abs*{(X, d)}, \forall r \in \R^{+}, N_{r}(p) \in \mathcal{O}(X)
  \]
\end{theorem}

\begin{proof}
  Suppose $q \in N_{r}(p)$. Then
  \[
    \begin{array}{c}
      d(p, q) < r \\
      \Downarrow  \\
      \veps \coloneqq r - d(p, q) > 0
    \end{array}
  \]
  Suppose $s \in N_{\veps}(q)$. Then
  \[
    \begin{array}{c}
      d(q, s) < \veps                                      \\
      \Downarrow                                           \\
      d(p, s) \leq d(p, q) + d(q, s) < d(p, q) + \veps = r \\
      \Downarrow                                           \\
      s \in N_{r}(p)                                       \\
    \end{array}
  \]
  Since $s \in N_{\veps}(q)$ was arbitrary,
  \[
    \begin{array}{c}
      N_{\veps}(q) \subseteq N_{r}(p) \\
      \Downarrow                      \\
      q \in \paren*{N_{r}(p)}^{\circ}
    \end{array}
  \]
\end{proof}

\begin{theorem} \label{an1:thm:limit-point-neighborhood-contains-infinite-elements}
  \[
    \forall E \subseteq \abs*{(X, d)}, \forall p \in E', \forall r \in \R^{+}, \abs*{E \cap N_{r}(p)} \geq \aleph_{0}
  \]
\end{theorem}

\begin{proof}
  Suppose $E \subseteq X$. Assume
  \[
    \exists p \in E', \exists r \in \R^{+}, \abs*{E \cap N_{r}(p)} < \aleph_{0}
  \]
  Define
  \[
    S \coloneqq E \cap \widetilde{N}_{r}(p) \neq \vnot
  \]
  Then
  \[
    \begin{array}{c}
      S \subseteq E \cap N_{r}(p)                                   \\
      \by{\Cref{an1:thm:subset-of-finite-set-is-finite}} \Downarrow \\
      \abs*{S} < \aleph_{0}
    \end{array}
  \]
  Define
  \[
    T \coloneqq \setbuilder*{d(p, q)}{q \in S} \neq \vnot
  \]
  Then
  \[
    \begin{array}{c}
      p \notin S \\
      \Downarrow \\
      T \subseteq \R^{+}
    \end{array}
  \]
  Define $f: S \to T$ by
  \[
    \forall q \in S, f(q) \coloneqq d(p, q)
  \]
  Then
  \[
    \begin{array}{c}
      \forall t \in T, \exists q \in S, t = d(p, q)                            \\
      \Downarrow                                                               \\
      f: S \twoheadrightarrow T                                                \\
      \by{\Cref{an1:thm:surjection-from-finite-set-implies-finite}} \Downarrow \\
      \abs*{T} < \aleph_{0}                                                    \\
      \by{\Cref{an1:thm:finite-ordered-set-has-extremum}} \Downarrow           \\
      \exists \veps \in \R^{+}, \veps = \min T
    \end{array}
  \]
  This implies
  \[
    \begin{array}{c}
      \forall q \in S, d(p, q) \geq \veps     \\
      \Downarrow                              \\
      E \cap \widetilde{N}_{\veps}(p) = \vnot \\
      \Downarrow                              \\
      \bot
    \end{array}
  \]
  Therefore,
  \[
    \abs*{E \cap N_{r}(p)} \geq \aleph_{0}
  \]
\end{proof}

\begin{theorem} \label{an1:thm:complement-of-open-set-is-closed}
  \[
    \forall E \subseteq \abs*{(X, d)}, E \in \mathcal{O}(X) \iff E^{c} \in \mathcal{C}(X)
  \]
\end{theorem}
